\subsection{Learning Controllers}\label{sec:learning}
Each animal owns a controller selected from Table~\ref{tab:algorithms}. Network
settings and learning-rate multipliers can differ by species. Policy MLPs support
configurable tanh hidden layers; the legacy Q network retains its small
single-hidden-layer contract. A deeper neural network is an architecture choice,
not an additional RL algorithm. Recurrent PPO uses GRU-style state
\citep{cho2014gru}; QR-DQN estimates return quantiles rather than one scalar Q
value per action. Adam~\citep{kingma2015adam} is used by the policy and QR-DQN
controllers; the legacy Q learner uses SGD.

\begin{table}[t]
\centering\small
\caption{Implemented controllers. All parameter/buffer ownership is per animal.
These are repository implementations and adaptations, not wrappers around the
cited authors' training libraries.}\label{tab:algorithms}
\begin{tabularx}{\linewidth}{@{}lYY@{}}\toprule
Controller & Update mechanism & Distinct state or information boundary\\\midrule
Q network / DQN & Replay, target network, Huber TD gradients, SGD & Epsilon-greedy actions and bounded replay\\
PPO & Clipped policy objective, GAE, critic and Adam & Private recent rollout; configurable epochs/minibatches\\
A2C & One actor--critic update per rollout & Private value estimates and advantages\\
REINFORCE & Discounted finite-episode policy gradient & Actor-only trajectory; baseline from earlier episodes\\
GRPO & Group-relative clipped actor updates & Matched trial worlds, frozen reference policy, no critic\\
Recurrent PPO & PPO with chronological recurrent gradients & GRU hidden state and sequence inputs\\
QR-DQN & Quantile-regression Double DQN and Adam & Return quantiles and target network\\
Maskable PPO & PPO over immediate feasible actions & Stored decision-time masks; consent remains unknown\\
MAPPO variant & PPO with centralized critic & Local actor; private actor/critic; species-mean training reward\\\bottomrule
\end{tabularx}
\end{table}

\subsubsection{On-policy and off-policy updates}
The Q learner draws minibatches from its own bounded replay and refreshes its
target at a configured update interval. Policy learners retain the observations
and decision data required by their objective. PPO uses clipped likelihood
ratios, GAE, entropy regularization, an unclipped mean-square critic objective,
gradient clipping and optional approximate-KL stopping. Advantages are normalized
within an individual's rollout when configured. A2C uses the same actor--critic
infrastructure without repeated PPO epochs. REINFORCE consumes a completed finite
trajectory and does not bootstrap a time-limit value. These distinctions are
preserved when a death or episode boundary flushes a partial buffer.

The experimental MAPPO implementation combines local actor observations with a
versioned centralized critic context and species-mean training reward. It does
not replace all animals with a shared population policy, nor should it be
interpreted as an exact reproduction of a published benchmark implementation.
Action masking similarly reflects the current local rules, not perfect
knowledge of simultaneous movements, attacks or another animal's consent.

\subsubsection{Matched simulator GRPO trials}
A due focal animal samples $G$ alternative trajectories from the same current
world state. Nonfocal policies remain frozen; world and nonfocal action RNGs
restart from matched states, while focal action samples differ. Trial rewards
come from the ordinary physical rules. No trial calls a real controller's
learning/action methods, advances the real world, or records hypothetical kills
and births as real events.

Let $R_g$ be the discounted return of trial $g$. The implementation standardizes
returns within the group to obtain $A_g$, then applies a clipped policy-ratio
objective with a reference-policy KL term. Each trajectory contributes equally
to the group mean rather than being implicitly weighted by its length. Groups
with equal returns are discarded, including the KL update. This makes reward
spread an important diagnostic but does not ensure useful exploration.

Trials run sequentially with isolated mutable state; read-only geography can be
shared. If at most $U$ groups are scheduled per real step with horizon $H$,
the extra physical-step budget is bounded by $UGH$, before early termination.
Group size, horizon, trial frequency and per-step scheduling must therefore be
reported when comparing algorithms. Matching real-world steps alone does not
match GRPO's computational or simulated-experience budget.
